\chapter*{Agradecimientos}
%\addcontentsline{toc}{chapter}{Agradecimientos} % si queremos que aparezca en el índice
\markboth{AGRADECIMIENTOS}{AGRADECIMIENTOS} % encabezado 

Puede que haya personas que deberían aparecer en esta página y no estén, pues para agradecer a todos los que me han ayudado, apoyado, y acompañado durante estos años, necesitaría una memoria entera.

Ahora bien, durante el tiempo que he estado inmersa, no solo en este proyecto, sino en la carrera, ha habido momentos buenos, también otros que quizás no lo fueron tanto, pero algo que siempre se mantuvo ahí, presente, pasara lo que pasase, fue el apoyo incondicional de mi madre, Fátima. A ella le quiero agradecer en primer lugar, pues ha sido mi pilar, mi apoyo y \textit{el viento que quita las nubes de mi sol nublado.} Ella ha estado presente en todas las noches en vela estudiando, tratando de que funcionaran simulaciones, y en todos los momentos de alegría cuando conseguía algo que me hacía sentir orgullosa.

Tampoco puedo dejar sin agradecer a mi familia, que han estado ahí, animándome siempre, quitando hierro a problemas que parecían un túnel sin salida, y celebrando mis logros como si fueran suyos. Gracias a mi padre, a mi yaya, a mis abuelos, a todos mis tíos y tías (Michel, Rebeca, Chus, Jorgito...), y a mis primos, tanto mayores como pequeños, que me han acompañado en esta etapa de mi vida y que han hecho que todo sea más fácil y llevadero.

Además, quiero nombrar a mis amigos, Alberto, Andrea, Ariadna e Inés, por estar ahí siempre, por salir a dar una vuelta en mis momentos de desesperación, por hacerme reír cuando más lo necesitaba, por escucharme, y por darme consejos que me han ayudado a crecer.

Además, durante este tiempo he tenido la suerte de conocer a muchas personas que me han ayudado a llegar a ser quién soy y a aprender.

En primer lugar, (albano y solano)

Durante este camino también he conocido a numerosos compañeros que han hecho que las largas jornadas en la universidad fueran más llevaderas, no solo con los momentos de alegría sino compartiendo la misma desesperación cuando algo no iba bien. Quiero dar las gracias a mis amigos de aeroespacial y mi doble grado: mi Jorgito, Silvia, Ana, María, María F., Manu, Alejandro, Paco, Andy, Gabriela, Jorge... Y muchos más que no puedo poner aquí. Espero seguir compartiendo momentos con vosotros toda la vida, os voy a llevar conmigo siempre.

Además, quiero agradecer a mis amigos de teleco, Sergio, David, Pedrito e Íker que han estado ahí desde el principio también y quienes estos últimos años se han convertido en una parte muy importante de mí, ocupando un lugar en mi corazón. Gracias por todo vuestro apoyo.

Por último gracias a todos esos amigos que son familia, vosotros ya sabéis quienes sois.
